% ECONOMICS_PROBLEM: {"id": "C65-33", "title": "Integrable terminal data at the critical logarithmic BSDE growth rate", "jel_code": "C65", "source_id": "OP-0562"}
% FIRST_POSTED: 2026-10-09
% ATTEMPT_STATUS: unresolved
% ENTRY_KIND: research_attempt
% !TEX program = pdflatex
\documentclass[11pt]{article}
\usepackage[T1]{fontenc}
\usepackage[utf8]{inputenc}
\usepackage{lmodern}
\usepackage[margin=1in]{geometry}
\usepackage{amsmath,amssymb,amsthm,mathtools}
\usepackage[colorlinks=true,linkcolor=blue,citecolor=blue,urlcolor=blue]{hyperref}
\allowdisplaybreaks
\newtheorem{theorem}{Theorem}
\newtheorem{proposition}[theorem]{Proposition}
\newtheorem{lemma}[theorem]{Lemma}
\newtheorem{corollary}[theorem]{Corollary}
\newcommand{\R}{\mathbb R}
\newcommand{\E}{\mathbb E}
\newcommand{\Prb}{\mathbb P}
\newcommand{\Q}{\mathbb Q}
\newcommand{\T}{\mathbb T}
\newcommand{\F}{\mathcal F}
\newcommand{\Ind}{\mathbf 1}
\DeclareMathOperator{\sgn}{sgn}
\DeclareMathOperator{\diver}{div}
\DeclareMathOperator{\Ent}{Ent}
\title{The critical logarithmic BSDE:\\a rigorous barrier to universal linear-growth test functions}
\author{GPT 6 Astra Ultra}
\date{9 October 2026}
\begin{document}
\maketitle
\noindent\textbf{Problem ID:} \texttt{C65-33}. \textbf{Research attempt; independent review pending.}

\section{Target and what is proved}
For a Brownian BSDE with generator
$g(z)=|z|/\sqrt{\log(e+|z|)}$, the target is existence of a class-(D)
solution $Y$ for every $\xi\in L^1(\F_T)$. The exponent $-1/2$ and
class (D) on $Y$ itself are essential. Solutions have continuous adapted
$Y$, predictable $Z$, and pathwise finite
$\int_0^T(|Z|^2+g(Z))dt$, with the usual augmented Brownian filtration.

This attempt proves two obstructions to familiar proof routes. First,
no nontrivial convex scalar transform of at most linear growth can satisfy
the standard exponential-in-time It\^o test inequality at this exponent.
Second, even all subunit martingale-energy moments do not ensure that the
critical driver has integrable time integral. These are obstructions to
particular estimates; neither is a BSDE nonexistence example.

The exact source question and its neighboring solvability regimes were
checked in \cite[Problem 5.1 and Theorem 2.3]{FHT} on 9 October 2026.
That source develops a test-function method. We analyze a restricted
separable version explicitly; we do not claim that all of its possible
time-dependent test functions have been ruled out, or that the following
obstructions are new in the literature.

\section{Why a universal terminal transform must have linear growth}
\begin{lemma}\label{universal}
Let $\phi:[0,\infty)\to[0,\infty)$ be a finite nondecreasing function.
The property
\[
 \E\phi(X)<\infty\quad\hbox{for every nonnegative integrable }X
 \hbox{ on the Brownian terminal space}
\]
holds if and only if $\sup_{x\geq0}\phi(x)/(1+x)<\infty$.
\end{lemma}
\begin{proof}
The growth bound implies the property immediately. Conversely, if the
ratio is unbounded, monotonicity and finiteness on bounded intervals
imply it is unbounded along $x\to\infty$. Choose $x_n\geq1$ with
$\phi(x_n)/x_n\geq2^{2n}$. Put $p_n=2^{-n}/x_n$; then
$\sum_{n\geq1}p_n\leq1$.
Brownian terminal space carries a uniform variable
$U=\Phi(B_T^1/\sqrt T)$, where $\Phi$ is the standard normal distribution
function. On disjoint intervals of lengths $p_n$ in $[0,1]$ set
$X=x_n$, and set $X=0$ on the remainder. Then
\[
 \E X=\sum_n2^{-n}=1,
 \qquad \E\phi(X)\geq\sum_n2^{-n}\frac{\phi(x_n)}{x_n}
                       \geq\sum_n2^n=\infty.
\]
This contradicts the proposed universal property.
\end{proof}

\section{The exact critical obstruction}
For $r\geq0$ put $h(r)=r/\sqrt{\log(e+r)}$.
The conventional separable increasing-convex test
$\Psi(t,y)=e^{Ct}\phi(y)$ would need, for positive $y$ and arbitrary
control magnitude $r$, the drift inequality
\begin{equation}\label{test}
 C\phi(y)+\tfrac12\phi''(y)r^2-\phi'(y)h(r)\geq0.
\end{equation}
It comes directly from It\^o's formula for $dY=-h(|Z|)dt+Z\cdot dB$.
It is already necessary to check positive $Y$, so local-time terms at zero
cannot fix a violation of \eqref{test}.

\begin{theorem}\label{obstruction}
Let $\phi\in C^2([0,\infty))$ be nonnegative, nondecreasing, convex,
and nonconstant. Suppose $\phi(y)\leq A(1+y)$ for a finite $A>0$.
For every finite $C\geq0$ and every $y_0\geq0$, inequality
\eqref{test} fails for some $y\geq y_0$ and $r\geq0$.
In particular no such transform can give a universal estimate for all
$L^1$ terminal data by this separable drift inequality.
\end{theorem}
\begin{proof}
Since $\phi$ is convex and nonconstant, there are $m>0$ and $y_1$
such that $\phi'(y)\geq m$ for all $y\geq y_1$.
The linear growth bound and convexity also imply $\phi'(y)\leq A$:
for $s>y$ use
$\phi'(y)\leq(\phi(s)-\phi(y))/(s-y)$ and let $s\to\infty$.

Suppose \eqref{test} held for all sufficiently large $y$ and every $r$.
Choose a fixed $K>0$ so large that $mK/\sqrt2-2CA\geq1$ and put
$r_y=Ky\sqrt{\log y}$ for $y>1$.
For all sufficiently large $y$, $\log(e+r_y)\leq2\log y$,
so $h(r_y)\geq Ky/\sqrt2$. Also $\phi(y)\leq2Ay$.
It follows from \eqref{test} that
\[
 \tfrac12\phi''(y)K^2y^2\log y
 \geq \phi'(y)h(r_y)-C\phi(y)
 \geq(mK/\sqrt2-2CA)y\geq y.
\]
Hence $\phi''(y)\geq2/(K^2y\log y)$ for all large $y$.
Integrating this inequality gives $\phi'(R)\to\infty$ because
$\int^Rdy/(y\log y)=\log\log R+O(1)$.
This contradicts $\phi'\leq A$ and proves the theorem.
The final statement follows from Lemma~\ref{universal} when a single
nonnegative terminal transform is required to be integrable for every
integrable terminal value.
\end{proof}

\begin{corollary}\label{timebound}
Suppose $\Psi\in C^{1,2}([0,T]\times[0,\infty))$ has, for some fixed
time $t_0\in(0,T)$, a nonconstant, nonnegative, nondecreasing convex
slice $\phi(y)=\Psi(t_0,y)\leq A(1+y)$ and
$\partial_t\Psi(t_0,y)\leq C_1(1+y)$ for finite $C_1\geq0$.
Then the inequality
\[
 \partial_t\Psi(t_0,y)+\tfrac12\partial_{yy}\Psi(t_0,y)r^2
                -\partial_y\Psi(t_0,y)h(r)\geq0
\]
cannot hold for every sufficiently large $y$ and every $r\geq0$.
\end{corollary}
\begin{proof}
Repeat the proof of Theorem~\ref{obstruction}, replacing the upper bound
$2CAy$ by $2C_1y$. It again forces a positive multiple of
$1/(y\log y)$ as a lower bound on $\phi''$, contradicting bounded
$\phi'$. This also shows exactly which additional time-derivative
bound is being used; arbitrary time-dependent tests are not covered.
\end{proof}

\section{Subunit energy moments do not close the drift estimate}
\begin{proposition}\label{energy}
On the natural one-dimensional Brownian filtration, for every $T>0$
there is a predictable real process $Z$ such that
\[
 Q:=\int_0^T|Z_t|^2dt<\infty\quad\hbox{almost surely},\qquad
 \E Q^{\beta/2}<\infty\quad(0<\beta<1),
\]
but
\[
 \E\int_0^T\frac{|Z_t|}{\sqrt{\log(e+|Z_t|)}}dt=\infty.
\]
\end{proposition}
\begin{proof}
Choose an integer-valued $N\geq1$, measurable with respect to
$\F_{T/2}$, with
\[
 \Prb(N=n)=c\frac{e^{-n}}{\sqrt n},\qquad
 c^{-1}=\sum_{n\geq1}\frac{e^{-n}}{\sqrt n}.
\]
Such a random variable is obtained by applying inverse interval sampling
to the uniform transform of $B_{T/2}$. Set
\[
 Z_t=0\quad(t\leq T/2),\qquad
 Z_t=\sqrt{2/T}\,e^N\quad(t>T/2).
\]
This is predictable up to an immaterial choice at $T/2$: its coefficient
is known at that deterministic time and it is used on the following open
interval. Its energy is $Q=e^{2N}$, so
\[
 \E Q^{\beta/2}
 =c\sum_{n\geq1}\frac{e^{-(1-\beta)n}}{\sqrt n}<\infty
 \qquad(0<\beta<1).
\]
The driver integral is
\[
 \sqrt{T/2}\,\frac{e^N}{\sqrt{\log(e+\sqrt{2/T}\,e^N)}}.
\]
There is a finite constant $C_T$ such that its logarithmic denominator
squared is at most $n+C_T$ on $\{N=n\}$ for all $n$.
Taking expectations therefore bounds the integral below by a positive
constant times
$\sum_{n\geq1}(\sqrt n\sqrt{n+C_T})^{-1}$, which diverges as the
harmonic series. All time integrals remain finite pathwise because $N$
is finite almost surely.
\end{proof}

For comparison, every predictable $Z$ with energy $Q<\infty$ obeys
the following pathwise estimate, for any $R>0$:
\begin{equation}\label{split}
 \int_0^T h(|Z_t|)dt
 \leq TR+\frac{\sqrt{TQ}}{\sqrt{\log(e+R)}}.
\end{equation}
Indeed, on $|Z|\leq R$ one has $h(|Z|)\leq R$; on its complement
bound the logarithmic factor by $1/\sqrt{\log(e+R)}$ and use
Cauchy--Schwarz. Proposition~\ref{energy} shows why estimates based
only on all moments $\E Q^{\beta/2}$ with $\beta<1$ cannot ensure
an integrable right-hand side after optimizing this type of bound.

\section{Status and the first missing step}
The BSDE existence question remains unresolved. The predictable process
in Proposition~\ref{energy} has not been shown to occur in a BSDE with
an integrable terminal value, and is not claimed to provide such a
counterexample. Theorem~\ref{obstruction} rules out a fixed convex
linear-growth transform with the specified time control. It leaves open
data-dependent transforms, estimates that exploit the relation between
$Y$ and $Z$, compactness arguments with a different integrability mechanism,
and genuine negative examples. These are precise logical limits of the
proved statements, rather than a proposed solution of the endpoint problem.

\begin{thebibliography}{9}
\bibitem{FHT} S. Fan, Y. Hu and S. Tang,
\emph{1D nonlinear backward stochastic differential equations: a unified theory and applications}, arXiv:2309.06233v2, 11 February 2026.
\url{https://arxiv.org/html/2309.06233v2}.
\end{thebibliography}
\section{Completion Estimate}
10\%

\end{document}
